\documentclass[10pt,a4paper]{amsart}
\usepackage{amsmath,amssymb,amsthm,graphicx}
% For commutative diagrams we use Xy-pic
\usepackage[all]{xy}

% We use 2cell for 2-commutative diagrams.
\xyoption{2cell}
% Theorem environments.
%
\theoremstyle{plain}
\newtheorem{theorem}[subsection]{Theorem}
\newtheorem{proposition}[subsection]{Proposition}
\newtheorem{lemma}[subsection]{Lemma}

\theoremstyle{definition}
\newtheorem{definition}[subsection]{Definition}
\newtheorem{example}[subsection]{Example}
\newtheorem{exercise}[subsection]{Exercise}
\newtheorem{situation}[subsection]{Situation}

\theoremstyle{remark}
\newtheorem{remark}[subsection]{Remark}
\newtheorem{remarks}[subsection]{Remarks}

\numberwithin{equation}{subsection}

% Macros
%
\def\lim{\mathop{\mathrm{lim}}\nolimits}
\def\colim{\mathop{\mathrm{colim}}\nolimits}
\def\Spec{\mathop{\mathrm{Spec}}}
\def\Hom{\mathop{\mathrm{Hom}}\nolimits}
\def\Ext{\mathop{\mathrm{Ext}}\nolimits}
\def\SheafHom{\mathop{\mathcal{H}\!\mathit{om}}\nolimits}
\def\SheafExt{\mathop{\mathcal{E}\!\mathit{xt}}\nolimits}
\def\Sch{\mathit{Sch}}
\def\Mor{\mathop{\mathrm{Mor}}\nolimits}
\def\Ob{\mathop{\mathrm{Ob}}\nolimits}
\def\Sh{\mathop{\mathit{Sh}}\nolimits}
\def\NL{\mathop{N\!L}\nolimits}
\def\CH{\mathop{\mathrm{CH}}\nolimits}
\def\proetale{{pro\text{-}\acute{e}tale}}
\def\etale{{\acute{e}tale}}
\def\QCoh{\mathit{QCoh}}
\def\Ker{\mathop{\mathrm{Ker}}}
\def\Im{\mathop{\mathrm{Im}}}
\def\Coker{\mathop{\mathrm{Coker}}}
\def\Coim{\mathop{\mathrm{Coim}}}

% Boxtimes
%
\DeclareMathSymbol{\boxtimes}{\mathbin}{AMSa}{"02}

%
% Macros for moduli stacks/spaces
%
\def\QCohstack{\mathcal{QC}\!\mathit{oh}}
\def\Cohstack{\mathcal{C}\!\mathit{oh}}
\def\Spacesstack{\mathcal{S}\!\mathit{paces}}
\def\Quotfunctor{\mathrm{Quot}}
\def\Hilbfunctor{\mathrm{Hilb}}
\def\Curvesstack{\mathcal{C}\!\mathit{urves}}
\def\Polarizedstack{\mathcal{P}\!\mathit{olarized}}
\def\Complexesstack{\mathcal{C}\!\mathit{omplexes}}
% \Pic is the operator that assigns to X its picard group, usage \Pic(X)
% \Picardstack_{X/B} denotes the Picard stack of X over B
% \Picardfunctor_{X/B} denotes the Picard functor of X over B
\def\Pic{\mathop{\mathrm{Pic}}\nolimits}
\def\Picardstack{\mathcal{P}\!\mathit{ic}}
\def\Picardfunctor{\mathrm{Pic}}
\def\Deformationcategory{\mathcal{D}\!\mathit{ef}}
\usepackage{cleveref}
\crefname{originallemma}{original lemma}{original lemmas}
\Crefname{originallemma}{Original lemma}{Original lemmas}
\makeatletter
\newlabel{algebra-lemma-local-dimension-zero-henselian}{{06RS}{source}{}{sourcepointer-1}{}}
\newlabel{algebra-lemma-local-dimension-zero-henselian@cref}{{[originallemma][0][]06RS}{[1][0][]source}{}{}{}}
\makeatother
\crefname{originallemma}{original lemma}{original lemmas}
\Crefname{originallemma}{Original lemma}{Original lemmas}
\makeatletter
\newlabel{algebra-lemma-ring-with-only-minimal-primes}{{04MG}{source}{}{sourcepointer-2}{}}
\newlabel{algebra-lemma-ring-with-only-minimal-primes@cref}{{[originallemma][0][]04MG}{[1][0][]source}{}{}{}}
\makeatother
\crefname{originallemma}{original lemma}{original lemmas}
\Crefname{originallemma}{Original lemma}{Original lemmas}
\makeatletter
\newlabel{cohomology-lemma-vanishing-for-profinite}{{0A3F}{source}{}{sourcepointer-3}{}}
\newlabel{cohomology-lemma-vanishing-for-profinite@cref}{{[originallemma][0][]0A3F}{[1][0][]source}{}{}{}}
\makeatother
\crefname{originalreference}{original reference}{original references}
\Crefname{originalreference}{Original reference}{Original references}
\makeatletter
\newlabel{lemma-local-rings-strictly-henselian}{{09AX}{source}{}{sourcepointer-4}{}}
\newlabel{lemma-local-rings-strictly-henselian@cref}{{[originalreference][0][]09AX}{[1][0][]source}{}{}{}}
\makeatother
\title{On an affine scheme whose points are closed with separably algebraically closed residue fields all higher cohomology vanishes}
\author{The Stacks Project Authors}
\date{stacks-project:etale-cohomology.tex:local-snapshot-20261011; GNU FDL 1.2 or later}
\begin{document}
\maketitle
\noindent\textit{Editorial note.} All delivered mathematical text is the author's own, copied from the pinned Stacks Project chapter file etale-cohomology.tex; no statement or proof step was rewritten, reordered or removed.\par\medskip
\noindent\textit{Editorial note.} The counted proof is the author's complete argument: the proof is written out in full: the profiniteness of the underlying space, the vanishing of topological cohomology, the strict henselianity of the local rings, and the comparison of the two cohomology theories.\par\medskip
\noindent\textit{Editorial note.} The label lemma-affine-only-closed-points names the counted statement, and the proof environment immediately following it in the source is delivered verbatim as its proof.\par\medskip
\noindent\textit{Editorial note.} The counted statement is carried by the author's own lemma environment in the chapter source.\par\medskip
\noindent\textit{Editorial note.} The statement and proof are copied from the pinned chapter file; the four cited inputs are kept as external pointers with their canonical Stacks tags.\par\medskip
\noindent\textit{Allowed standard prerequisites (editorial).}\begin{itemize}
\item The underlying space is profinite. Algebra, the lemma stating that a ring whose primes are all maximal has a profinite spectrum. Supplies that the underlying topological space of S is profinite.
\item Vanishing of cohomology on a profinite space. Cohomology, the lemma that the higher cohomology of an abelian sheaf on a profinite topological space vanishes. Supplies the triviality of the Zariski cohomology.
\item Local rings of such a scheme are strictly henselian. Algebra, the lemma that a local ring of dimension zero with separably algebraically closed residue field is strictly henselian, together with Same chapter, the comparison of etale and Zariski cohomology for schemes with strictly henselian local rings. Supplies the identification of etale with Zariski cohomology.
\end{itemize}\medskip
\ifdefined\hypertarget\hypertarget{sourcepointer-1}{}\fi
\noindent\textit{Source reference.} Original lemma 06RS: sources/stacks-chow/tags.txt, source lines 8651--8651. Prerequisite pointer invoked at proof line(s) 14296; it names a standard fact of a sibling chapter that the counted argument uses, and every step of the counted proof built on it is retained verbatim in the delivered body.\par\medskip
\ifdefined\hypertarget\hypertarget{sourcepointer-2}{}\fi
\noindent\textit{Source reference.} Original lemma 04MG: sources/stacks-chow/tags.txt, source lines 6042--6042. Prerequisite pointer invoked at proof line(s) 14291; it names a standard fact of a sibling chapter that the counted argument uses, and every step of the counted proof built on it is retained verbatim in the delivered body.\par\medskip
\ifdefined\hypertarget\hypertarget{sourcepointer-3}{}\fi
\noindent\textit{Source reference.} Original lemma 0A3F: sources/stacks-chow/tags.txt, source lines 12738--12738. Prerequisite pointer invoked at proof line(s) 14294; it names a standard fact of a sibling chapter that the counted argument uses, and every step of the counted proof built on it is retained verbatim in the delivered body.\par\medskip
\ifdefined\hypertarget\hypertarget{sourcepointer-4}{}\fi
\noindent\textit{Source reference.} Original reference 09AX: sources/stacks-etale-cohomology/pinned-original-source.tex, source lines 14205--14205. Prerequisite pointer invoked at proof line(s) 14298; it names a standard fact of the same chapter that the counted argument uses, and every step of the counted proof built on it is retained verbatim in the delivered body.\par\medskip
\begin{lemma}
\label{lemma-affine-only-closed-points}
Let $S$ be an affine scheme such that
(1) all points are closed, and (2) all residue fields are separably
algebraically closed. Then
for any abelian sheaf $\mathcal{F}$ on $S_\etale$ we have
$H^i(S_\etale, \mathcal{F}) = 0$ for $i > 0$.
\end{lemma}

\begin{proof}
Condition (1) implies that the underlying topological space
of $S$ is profinite, see
Algebra, Lemma \ref{algebra-lemma-ring-with-only-minimal-primes}.
Thus the higher cohomology groups of an abelian sheaf on the topological
space $S$ (i.e., Zariski cohomology) is trivial, see
Cohomology, Lemma \ref{cohomology-lemma-vanishing-for-profinite}.
The local rings are strictly henselian by
Algebra, Lemma \ref{algebra-lemma-local-dimension-zero-henselian}.
Thus \'etale cohomology of $S$ is computed by Zariski cohomology
by Lemma \ref{lemma-local-rings-strictly-henselian}
and the proof is done.
\end{proof}

\end{document}
